\documentclass[11pt,a4paper]{article}
\usepackage[T1]{fontenc}
\usepackage[utf8]{inputenc}
\usepackage[english]{babel}
\usepackage{amsmath,amssymb}
\usepackage[margin=2cm]{geometry}
\setlength{\parindent}{0pt}
\setlength{\emergencystretch}{3em}
\allowdisplaybreaks

\title{About \texttt{integral\_conventions.tex}}
\author{}
\date{}
\begin{document}
\maketitle
This note explains what the test data in \texttt{integral\_conventions.tex} is used for. That file is input for the engine in \texttt{main.c}; it is read by the test \texttt{test\_integrals.c} in this folder.

\begin{enumerate}
\item \textbf{Purpose.} The file holds 16 true statements, written by hand, that fix the conventions of the reader (paragraph 9j of \texttt{parser.c}) and of the verdict (chapter 30 of \texttt{main.c}). The values were computed by hand.

\item \textbf{What it exercises.}
\begin{itemize}
\item Where the differential may stand: at the end, $f \, dx$; in a numerator, $\frac{dx}{g}$ and $\frac{x \, dx}{g}$; closing a sum, $\int x^{2} + 1 \, dx$, which is the integral of $x^{2} + 1$; after an argument that starts with a number, $\int \sin 2x \, dx$; alone, $\int 0 \, dx = c$; with another letter, $dt$.
\item The constant of the antiderivative: $c$, $C$ or $k$, and a constant already inside $F$ ($x^{2} + 5 + k$).
\item The meaning of an antiderivative: on every open interval where both sides are defined. So $\ln x$ and $\ln(-x)$ are antiderivatives of $\frac{1}{x}$ as $\ln|x|$ is, each where it is defined, and $\ln|3x|$ differs from $\ln|x|$ by a constant. The verdict says when the antiderivative is not defined everywhere the integrand is.
\item Definite integrals: bounds in either order of writing, $\int_{a}^{b}$ and $\int^{b}_{a}$; reversed bounds, $\int_{1}^{0} = -\int_{0}^{1}$; $\pi$ and $e$ in the bounds and in the value.
\end{itemize}

\item \textbf{Expected result.} The line \texttt{PASS: 16 of 16 statements proved, each with a written proof}.

\item \textbf{How to run it.} From the top folder of the repository:
\begin{verbatim}
cc -std=c11 -O2 -Wall -Wextra -o /tmp/test_integrals \
    tests/integrals/test_integrals.c && /tmp/test_integrals
\end{verbatim}
\end{enumerate}
\end{document}
